% ECONOMICS_PROBLEM: {"id": "Q58-103", "title": "Clean-technology policy sequences", "jel_code": "Q58", "source_id": "OP-0439"}
% FIRST_POSTED: 2026-10-09
% ATTEMPT_STATUS: unresolved
% ENTRY_KIND: research_attempt
\documentclass[11pt,a4paper]{article}
\usepackage[T1]{fontenc}
\usepackage[utf8]{inputenc}
\usepackage[margin=27mm]{geometry}
\usepackage{amsmath,amssymb,amsthm,mathtools}
\usepackage[hidelinks]{hyperref}
\setlength{\parskip}{5pt}
\setlength{\parindent}{0pt}
\newtheorem{theorem}{Theorem}[section]
\newtheorem{proposition}[theorem]{Proposition}
\newtheorem{lemma}[theorem]{Lemma}
\newtheorem{corollary}[theorem]{Corollary}
\theoremstyle{remark}
\newtheorem{remark}[theorem]{Remark}
\newcommand{\R}{\mathbb R}
\newcommand{\E}{\mathbb E}
\newcommand{\Prob}{\mathbb P}
\newcommand{\ind}{\mathbf 1}
\DeclareMathOperator{\Var}{Var}
\DeclareMathOperator{\KL}{KL}
\DeclareMathOperator{\TV}{TV}
\DeclareMathOperator{\OPT}{OPT}
\title{An explicit threshold for innovation before carbon pricing}
\author{GPT 6 Astra Ultra}
\date{9 October 2026}
\begin{document}\maketitle
\textbf{Problem ID:} \texttt{Q58-103}. \textbf{Status:} Unresolved. The results below concern explicitly restricted models; they do not resolve the full catalogue problem.

\section{Question and a two-date restriction}
The catalogue asks how innovation, infrastructure, carbon taxes, transfers, and political support should be sequenced over a long transition with uncertain parameters. Here we solve a deterministic two-date restriction with a knowledge spillover, irreversible early research, an emissions requirement, and a fixed activation cost. We then characterize fiscal and political constraints through a convex program. This is a model result; it supplies no empirical learning coefficient or general thirty-year optimum.

Dates are $0,1$. Abatement is $a_0,a_1\ge0$, and cumulative required abatement is $D>0$. Early public research is $r\ge0$, installed irreversibly at date zero and giving knowledge $Z_1=r$. Research is usable only at date one. Abatement and research costs, in date-zero resource units, are
\[
 C(a_0,a_1,r)=\frac{c_0a_0^2}{2}+\frac{k r^2}{2}
 +\beta\left(\frac{c_1a_1^2}{2}-\vartheta r a_1\right),
 \tag{1}
\]
where $c_0,c_1,k>0$, $\beta\in(0,1]$, and $\vartheta\ge0$. The emissions constraint is $a_0+a_1\ge D$. Baseline emissions minus $D$ is the cumulative budget. Cost (1) represents a local quadratic technology; its domain and convexity conditions are part of this restriction, not empirical assertions about global technology.

If research is used, an additional fixed resource cost $f\ge0$ is paid. The government can choose regime 0, requiring $r=0$ and no activation cost, or regime 1, allowing $r\ge0$ and adding $f$. Assume
\[
 kc_1>\beta\vartheta^2,\qquad
 d_0=\beta c_1,\qquad
 d_1=\beta c_1-\frac{\beta^2\vartheta^2}{k}>0. \tag{2}
\]
The quadratic form is then strictly positive definite. A representative household has linear utility in remaining resources. Taxes and rebates are internal transfers, so minimizing real cost is equivalent to maximizing this restricted household welfare.

\section{Solving the sequencing choice}
\begin{theorem}[Exact innovation threshold]
In regime $j\in\{0,1\}$, without additional fiscal or political constraints, the unique cost-minimizing abatement levels are
\[
 a_{0j}=\frac{d_jD}{c_0+d_j},\qquad
 a_{1j}=\frac{c_0D}{c_0+d_j}.
\]
Research is $r_0=0$ in regime 0 and $r_1=\beta\vartheta a_{11}/k$ in regime 1. The corresponding minimized variable costs are
\[
 V_j(D)=\frac{D^2}{2}\frac{c_0d_j}{c_0+d_j}.
\]
Therefore activation is strictly optimal exactly when
\[
 f<\Delta(D):=\frac{D^2}{2}
 \frac{c_0^2(d_0-d_1)}{(c_0+d_0)(c_0+d_1)}. \tag{3}
\]
If $\vartheta>0$, activation lowers date-zero abatement and raises date-one abatement, even though innovation itself occurs early.
\end{theorem}
\begin{proof}
For fixed $a_1$, complete the square:
\[
 \frac{k r^2}{2}-\beta\vartheta r a_1
 =\frac{k}{2}\left(r-\frac{\beta\vartheta a_1}{k}\right)^2
 -\frac{\beta^2\vartheta^2a_1^2}{2k}.
\]
The minimizing $r$ is nonnegative, so regime 1 reduces to $c_0a_0^2/2+d_1a_1^2/2$. Regime 0 has the same form with $d_0$. Both coefficients are positive. At a minimum $a_0+a_1=D$: if the sum were larger, proportional contraction would reduce the positive quadratic objective. Substitute $a_1=D-a_0$ and differentiate the strictly convex quadratic to obtain the displayed levels and values. Subtract $V_1$ from $V_0$ to obtain (3). Strict decrease of $d_1$ relative to $d_0$ lowers $dD/(c_0+d)$ and raises its complementary share, proving the sequencing statement.
\end{proof}

The comparison conditions on the \emph{same} emissions requirement. Relaxing the budget in one regime would change the economic question. Nor does early research necessarily imply early abatement: the knowledge spillover changes the intertemporal marginal-cost ratio.

\begin{proposition}[Implementation by taxes and a research payment]
Suppose firms take public knowledge $r$ as given and independently minimize
\[
 c_0a_0^2/2-\tau_0a_0,\qquad
 c_1a_1^2/2-\vartheta r a_1-\tau_1a_1
\]
over nonnegative abatement. Every interior optimum in Theorem 2.1 is implemented by
\[
 \tau_0=c_0a_0,\qquad \tau_1=c_1a_1-\vartheta r.
\]
In the research regime $\tau_0=\beta\tau_1$ and $\tau_1>0$. A research contractor with cost $kr^2/2$ supplies the optimal $r$ when paid marginal subsidy $\beta\vartheta a_1$ per unit of research, taking its announced value as given.
\end{proposition}
\begin{proof}
The two private objectives are strictly convex, so their first-order conditions uniquely give the stated abatement choices. At the optimum $c_0a_0=d_1a_1=\beta(c_1a_1-\vartheta r)$; positivity follows from (2). The contractor maximizes a linear payment times $r$ minus $kr^2/2$, whose unique maximizer is $\beta\vartheta a_1/k$.
\end{proof}

This is a commitment implementation, with taxes on baseline emissions equivalent to the displayed abatement incentives up to a fixed term. The research payment generally costs more fiscally than the contractor's real resource cost. That distinction matters once public spending is constrained.

\section{Political support, fiscal limits, and a complete certificate}
To make this distinction explicit, use direct public procurement of research at real cost $kr^2/2$; the contractor-subsidy implementation above is a separate unrestricted benchmark. Let $t\ge0$ be a date-zero transfer to households burdened by adjustment. Assume support is the known index
\[
 S=s_{\rm base}+t-\chi a_0,\qquad \chi\ge0,
\]
and that a dollar of transfer has marginal public-finance resource cost $\lambda\ge0$ beyond its transfer value. Support must exceed $s_0$, and research plus transfers must obey a fiscal limit $F$. In a fixed activation regime $j$, add constraints
\[
 t\ge s_0-s_{\rm base}+\chi a_0,\qquad
 kr^2/2+t+jf\le F,
\]
and minimize $C+\lambda t+jf$ subject also to the emissions requirement and nonnegativity. In regime 0 impose $r=0$.

\begin{proposition}[Feasibility and constrained optimality]
In regime 1 this is a convex optimization problem under (2). If its feasible set is nonempty, a minimum is attained. If there is a point satisfying all inequality constraints strictly, a feasible point is optimal if and only if it has nonnegative multipliers $\eta,\zeta,\xi$ and $\rho_0,\rho_1,\rho_r,\rho_t$ satisfying complementarity and
\begin{align*}
 c_0a_0-\eta+\chi\zeta-\rho_0&=0,\\
 \beta c_1a_1-\beta\vartheta r-\eta-\rho_1&=0,\\
 kr-\beta\vartheta a_1+\xi kr-\rho_r&=0,\\
 \lambda-\zeta+\xi-\rho_t&=0.
\end{align*}
The multipliers $\eta,\zeta,\xi$ correspond respectively to $D-a_0-a_1\le0$, $s_0-s_{\rm base}+\chi a_0-t\le0$, and $kr^2/2+t+f-F\le0$; the $\rho$ multipliers correspond to nonnegativity. Complementarity means each multiplier times its own constraint residual is zero.
\end{proposition}
\begin{proof}
The Hessian of $C$ has positive $c_0$ block and $(r,a_1)$ block determinant $\beta k c_1-\beta^2\vartheta^2>0$, so $C$ is strictly convex in these three variables. Every constraint defines a closed convex set. Along a minimizing sequence with bounded objective, positive definiteness bounds $(a_0,a_1,r)$, while the fiscal constraint bounds $0\le t\le F-f$. A convergent subsequence is feasible and attains the minimum.

Form the Lagrangian with the specified signs. Differentiation yields the equations. If the equations and complementarity hold, convexity gives that the candidate globally minimizes the Lagrangian; feasibility and nonnegative multipliers then give its objective no greater than that of any feasible competitor. Conversely, strict feasibility permits separation of the convex set of achievable strict objective improvement and constraint residuals from the negative orthant. Normalize the separating objective coefficient to one; strict feasibility rules out a zero coefficient. The remaining coefficients are nonnegative multipliers. The supporting condition gives Lagrangian minimization and complementarity, and differentiability gives the equations.
\end{proof}

The theorem compares two constrained regime values after feasibility is checked; (3) need not survive binding fiscal or political constraints. The expectation constraints of the full catalogue become deterministic ones here. Uncertain learning, entry, household heterogeneity, and changing political coalitions are not hidden inside these coefficients.

\section{Evidence still needed}
The model predicts a threshold in the fixed cost and in the required cumulative reduction, conditional on an identified $\vartheta$. Recovering that coefficient requires variation in early research or knowledge that is separated from anticipated future demand and concurrent infrastructure. Estimating $\chi$ requires a declared measurement model for support. A calibration that selects these coefficients solely to make (3) hold would not answer the empirical agenda. Stochastic nonanticipating policies, distributional welfare weights, and the full 31-date transition remain unresolved.

\begin{thebibliography}{9}
\bibitem{Stern} N. Stern (2022), \emph{A Time for Action on Climate Change and a Time for Change in Economics}, Economic Journal 132(644), 1259--1289. DOI: \href{https://doi.org/10.1093/ej/ueac005}{10.1093/ej/ueac005}. Primary author-institution record: \url{https://www.lse.ac.uk/granthaminstitute/publication/a-time-for-action-on-climate-change-and-a-time-for-change-in-economics-2/}.
\bibitem{BS} A. Bilal and J. H. Stock, \emph{A Guide to Macroeconomics and Climate Change}, NBER Working Paper 33567, March 2025, revised January 2026 according to the primary record. \url{https://www.nber.org/papers/w33567}. These sources motivate dynamic policy and empirical calibration; neither is claimed to contain the quadratic threshold model proved here.
\end{thebibliography}

\end{document}
